\chapter{项目背景与机会判断}

\section{政策环境与时代窗口}

智舆的提出依托两个清晰条件：政策端持续推动数字政府与智慧城市建设，技术端多模态分析与大模型能力进入可工程化阶段。河南省数字化转型与信阳本地治理平台升级，为基层治理场景引入 AI 辅助决策提供了现实入口。项目并非抽象概念拼接，而是对地方治理中长期存在的舆情识别与处置效率问题的针对性回应。

\begin{figure}[htbp]
	\t\centering
	\t\includegraphics[width=0.88\textwidth]{chapter1-policy-system.png}
	\t\caption{国家-省-市三级政策支持体系图}
\end{figure}

从政策支持看，国家层面明确推进数字政府与人工智能应用落地，省市层面推动治理能力现代化与数据资源整合。地方治理的制度与基础设施已经形成，使得“治理流程数字化+舆情智能分析”的结合具备现实土壤。

\section{传统舆情治理的核心痛点}

\begin{table}[htbp]
	\t\centering
	\t\caption{传统舆情治理主要痛点}
	\t\begin{tabularx}{\textwidth}{|L{3cm}YY|}
		\t\t\toprule
		\t\t\tblhead{痛点类别} & \tblhead{具体表现}              & \tblhead{对业务造成的影响}  \\
		\t\t\midrule
		\t\t\rowcolors{2}{tablealt}{white}
		\t\t数据覆盖不足         & 平台来源分散，图片、短视频、评论和热线数据难以统一纳入 & 风险发现滞后，容易漏掉早期弱信号    \\
		\t\t分析依赖人工         & 研判人员需要手工筛选、分类、汇总与写报         & 响应速度慢，人工成本高，结论一致性不足 \\
		\t\t缺少处置建议         & 系统能报警，但难以形成面向管理者的处置方案比较     & 监测后仍然要依赖经验判断，决策链路长  \\
		\t\t下沉市场门槛高        & 成熟系统价格高、部署重、服务半径集中在一二线城市    & 地市县客户难以承担，长期存在工具缺口  \\
		\t\t空间表达不足         & 传统报表偏列表化、平面化，难以结合城市空间态势理解问题 & 指挥调度和跨部门协同效率受限      \\
		\t\t\bottomrule
		\t\end{tabularx}
\end{table}

市场空间并非萎缩。公开研究显示，全球舆情监测系统市场在 2024 年约 23.15 亿美元，2031 年预计达到 40.50 亿美元，中国市场 2025 年预计约 72.4 亿元。对智舆而言，机会不在“市场有多大”，而在下沉场景长期缺乏价格可承受、部署可落地、表达可理解的工具体系。

\begin{figure}[htbp]
	\t\centering
	\t\includegraphics[width=0.86\textwidth]{chapter1-market-growth.png}
	\t\caption{中国舆情监测市场规模增长趋势图（2020-2030）}
\end{figure}

\section{信阳场景为何适合作为切入点}

信阳兼具文旅流量、基层治理与城市品牌传播等特征，舆情高频且类型多样。景区服务、民生热线与校园治理等场景需要更快的风险识别与闭环处置能力。对于智舆而言，信阳不是抽象样本，而是具备业务场景、试点可达性与本地资源协同优势的首批落地环境。

智舆在信阳的试点布局覆盖景区文旅、基层民生、校园治理与企业品牌四类场景，形成多场景协同验证的样板体系。景区文旅场景聚焦投诉热点与节假日波动管理，基层民生场景关注高频投诉主题与跨渠道重复诉求识别，校园治理场景强调热点事件预警与情绪变化监测，企业品牌场景服务本地企业声量与风险管理需求。多场景协同既验证系统的适配能力，也为后续区域复制积累可参考的实施样板。
